% ============================================================
% ch24.tex —— 第 24 章（选读）黎曼的藏宝图：$\zeta$ 函数与未解之谜
% ============================================================
\chapter{黎曼的藏宝图：\texorpdfstring{$\zeta$}{ζ} 函数与未解之谜}

\noindent\textbf{本章为观光章}（选读，难度 $\bigstar\bigstar\bigstar$）：不解决任何问题，但你会亲眼看见\textbf{数学地图上的空白}长什么样——以及为什么全世界最聪明的头脑们为一个空白排队排了一百六十多年，悬赏一百万美元。

\section{$\zeta$：素数的全息照片}

第 22 章的欧拉机器上还留着一个没拧到底的旋钮：把"倒数和"里的指数 $1$ 换成变量 $s$：
\[
\zeta(s) \;=\; \sum_{n=1}^{\infty}\frac{1}{n^s} \;=\; 1 + \frac1{2^s} + \frac1{3^s} + \frac1{4^s} + \cdots ,
\]
且欧拉乘积照搬（每个齿轮换成 $\tfrac{1}{1-p^{-s}}$）：
\[
\zeta(s) \;=\; \prod_{p}\frac{1}{1-p^{-s}} \qquad (s > 1).
\]
$\zeta$ 读"截塔"，希腊字母 Z——像 $x, y$ 一样，它是个\textbf{变量}：给一个 $s$，吐一个数。摸几个值找手感：
\begin{itemize}
\item $s = 2$：$\zeta(2) = 1 + \tfrac14 + \tfrac19+\cdots = \dfrac{\pi^2}{6} \approx 1.6449$（巴塞尔问题，欧拉 1735 年的成名战——\textbf{平方倒数和里冒出圆周率}，第 21 章埋的线在此引爆）；
\item $s = 4$：$\zeta(4) = \dfrac{\pi^4}{90} \approx 1.0823$；$s=6$：$\dfrac{\pi^6}{945}$——偶数处全是"$\pi$ 的幂乘有理数"；
\item $s = 3$：$\zeta(3) \approx 1.2021$（阿培里常数，1978 年才被证明是无理数——奇数处的规律至今无人参透）；
\item $s \to 1^+$：$\zeta(s) \to \infty$（调和级数，第 21 章）——\textbf{$s=1$ 是悬崖}。
\end{itemize}
这座桥的妙处：\textbf{左岸看所有整数（$\sum n^{-s}$），右岸只看素数（$\prod_p$），唯一分解是桥墩}。$s$ 是旋钮：拧大，级数飞快收敛、齿轮温顺；拧向 $1$，世界暴动。素数的一切统计信息（密度、间隔、分布）都编码在"$\zeta(s)$ 随 $s$ 变化的形状"里——所以说 $\zeta$ 是素数的\textbf{全息照片}：一小片底片，存着全部立体信息。

\section{零点：函数的溺水点}

1859 年，黎曼（32 岁，为入选柏林科学院提交论文，是他一生唯一一篇数论论文，\textbf{八页}）做了一件前无古人的事：把 $\zeta$ 从实数轴搬到\textbf{复平面}——横轴实部 $\sigma$、纵轴虚部 $t$，即 $s = \sigma + \mathrm{i}t$（数轴升维成地图，$\mathrm{i}$ 把"数"扩成"位置"）。

搬过去的第一个难关：级数 $\sum n^{-s}$ 只在 $s > 1$（实部大于 $1$）时有意义，$s \le 1$ 级数发散，函数"不存在"。黎曼用一套叫\textbf{解析延拓}的合法外推，给 $\zeta$ 在\textbf{除 $s=1$ 外的整张复平面}上发了户口——好比录音带只录了一段，但旋律的规律足以推算出整首曲子（外推不是瞎编：任何两条"解析"的延拓路径必得同一答案，唯一性是复分析的铁律）。

现在可以问那个致命的问题了：\textbf{$\zeta(s)$ 在哪些点等于零？}（"溺水点"——函数沉到零的位置。）这个问题为什么价值连城？回收第 23 章的线索：$\pi(x)$ 与 $\mathrm{li}(x)$ 的\textbf{误差由零点控制}。黎曼导出了"显式公式"：
\[
\pi(x) \;=\; \underbrace{\mathrm{li}(x)}_{\text{主旋律}} \;-\; \sum_{\rho}\underbrace{\mathrm{li}(x^{\rho})}_{\text{泛音，由零点 } \rho \text{ 奏出}} \;+\; (\text{小杂音}) ,
\]
素数的分布 $=$ 一条光滑主旋律 $+$ 一串由零点奏出的"波动泛音"——\textbf{零点在哪，素数的噪声就怎么抖}。零点离"$\sigma = 1$"这条线越远，泛音衰减越快，预报越准。素数定理（1896）证的正是"直线 $\sigma=1$ 上没有零点"；想要更准的预报（误差 $O(\sqrt x\ln x)$ 级），就得把零点推到最远可能的位置。

复平面地图侦察结果（黎曼已画出，后世全部证实）：
\begin{itemize}
\item \textbf{平凡零点}：$s = -2, -4, -6, \dots$（负偶数），规规矩矩排在地图西侧实轴上，性质温和，不影响大局；
\item \textbf{非平凡零点}：全部住在\textbf{临界带} $0 < \sigma < 1$（地图中央的竖条）里，且关于两条镜面对称——实轴（$\rho$ 与共轭 $\bar\rho$ 成对）和\textbf{临界线 $\sigma = \tfrac12$}（$\rho$ 与 $1-\rho$ 成对，这条对称来自 $\zeta$ 的"函数方程"，是黎曼论文的头条技术）。
\end{itemize}
计算机出场（零点可以逐个数值验证）：前十万亿（$10^{13}$）个非平凡零点，\textbf{全部精确落在临界线 $\sigma = \tfrac12$ 上}，无一离线。第一个零点在 $\tfrac12 + 14.134725\mathrm{i}$。

\begin{caixiang}[黎曼猜想，1859 年]
$\zeta$ 的全部非平凡零点都在临界线 $\mathrm{Re}(s) = \tfrac12$ 上。
\end{caixiang}

若它成立：所有泛音被压到理论最小振幅，素数定理的误差降到最优——\textbf{素数的分布"随机中带着最大限度的规整"}。它还是数论的"批发商"：数百条定理以"假设黎曼猜想成立"为前提（类比：整个楼层押在一根柱子上）；它对广义黎曼猜想（其他 $\zeta$ 类函数）的推广，是椭圆曲线、代数数论大量结果的地基（第 32 章的 BSD 猜想就是其中一支）。

\section{悬赏与现状}

现状清点（截至本书写作）：
\begin{itemize}
\item \textbf{数值}：验证了临界线上前 $10^{13}$ 个零点，零失误（Odlyzko 等人的计算；还验证了"临界带内前若干个零点全部在线上"）；
\item \textbf{理论下界}：哈代 1914 年证临界线上有\textbf{无穷多}零点；塞尔伯格 1942 年证"正比例"；莱文森 1974 年推到 $\tfrac13$；康瑞 1989 年推到 \textbf{$40\%$}——四十年过去，$40\%$ 仍是纪录；
\item \textbf{等价与推论}：黎曼猜想 $\iff$ 素数定理最优误差 $\iff$ $\sum_{n\le x}\mu(n) = O(x^{1/2+\varepsilon})$（$\mu$ 是莫比乌斯函数，"素数个数的带符号会计"）等几十种马甲；
\item \textbf{悬赏}：希尔伯特 1900 年 23 问题的第 8 号；克雷研究所 2000 年千禧年七大难题，每题一百万美元——黎曼猜想是唯一同时在这两张榜单上的。
\end{itemize}

传闻希尔伯特说过：若五百年后复活，第一个问题就问"黎曼猜想证了吗"。为什么全世界相信它？三路证据：数值的零失误；类比的胜利（函数域上的"平行版黎曼猜想"已被德尔涅 1974 年用代数几何证明——几何世界全对，数的世界凭什么错）；物理的暗合（蒙哥马利 1972 年发现零点间距的统计规律与随机厄米矩阵特征值惊人一致，恰好是管理重原子核能级的数学——\textbf{素数的噪声里，似乎藏着量子力学的琴弦}，这条线索至今无人参透）。当然，证据再多也不算证明——欧拉乘积下 $\sum 1/p$ 的"十万亿个对"换不来一行证明，数学家比谁都清楚。

\begin{lishi}
格奥尔格·弗雷德里希·伯恩哈德·黎曼（1826--1866），高斯的学生，一生发表论文不超过十篇，却每一篇都开辟一个领域（黎曼几何——第 28 章宇宙形状的语言；黎曼面——复分析的地基；黎曼积分——第 17 章切条的正式化；……）。1859 年那八页数论论文是他为竞选柏林科学院院士的"就职献礼"（按惯例向院系献礼），高斯死后他接任哥廷根数学教席。1866 年死于肺结核，39 岁。那八页纸里他说了一句话："零点很可能全部位于……"（\textit{von der Richtigkeit dieser Vermutung\"ubrigens ein weniger sicherer Beweis}——"当然，对这一猜想的正确性，我还欠一个不太确定的证明"）——一百六十多年，全人类都在替他补这个证明。
\end{lishi}

\begin{dongshou}
\begin{itemize}
  \yisi 手算 $\zeta(2)$ 的前六项：$1 + 0.25 + 0.1111 + 0.0625 + 0.04 + 0.0278 = 1.4914$，对比 $\tfrac{\pi^2}{6} = 1.6449$（还差 $9.3\%$）；再算 $\zeta(4)$ 前四项 $1 + 0.0625 + 0.0123 + 0.0039 = 1.0787$，对比 $\tfrac{\pi^4}{90} = 1.0823$（只差 $0.3\%$）——$s$ 越大收敛越飞快（分母是 $n^s$，指数加速）。
  \yier 验证"大 $s$ 的 $\zeta$ 极限"：$\zeta(10) = 1 + \tfrac1{1024} + \tfrac1{59049} + \cdots \approx 1.000994$——$s$ 稍大，$\zeta$ 就贴着 $1$。这解释了欧拉乘积在 $s$ 大时为什么是"绝对收敛的老实人"，而 $s = 1$ 恰是暴动临界点。
  \yier 把 $\zeta(2)$ 的欧拉乘积写出来并数值逼近：$\prod_p (1-p^{-2})^{-1} = \tfrac{\pi^2}6$。注意齿轮的样子：素数 $p$ 的齿轮是 $\tfrac{1}{1-p^{-2}}$，比如 $p=2$ 是 $\tfrac43$、$p=3$ 是 $\tfrac98$、$p=5$ 是 $\tfrac{25}{24}$。前四个齿轮：$\tfrac43\times\tfrac98\times\tfrac{25}{24}\times\tfrac{49}{48} = \tfrac{44100}{27648} \approx 1.595$；再加 $11, 13, 17, 19$ 四个齿轮（各约 $1.008, 1.006, 1.003, 1.003$），八个齿轮 $\approx 1.628$——从下方稳稳爬向 $1.6449$。每个素数齿轮都在给 $\pi$ 投票。\textbf{这题请务必亲手乘一遍}：那是欧拉 1735 年震动欧洲的心跳。
\end{itemize}
\end{dongshou}

\begin{shaonao}
一个只有初等工具就能感受到的涟漪：$\zeta(-1) = 1 + 2 + 3 + 4 + \cdots$ 按求和定义当然发散，但解析延拓后黎曼算出
\[
\zeta(-1) = -\frac{1}{12} \qquad\Longleftrightarrow\qquad ``1+2+3+4+\cdots = -\tfrac1{12}"
\]
——网络疯传的"荒谬等式"。它不是普通加法（普通加法下发散），而是"把 $\zeta$ 合法延拓到 $-1$ 后的取值"。可以亲手体验"延拓"味道的是它的简化版——\textbf{阿贝尔求和}：给发散级数戴上"衰减因子 $x^n$"，在 $x<1$ 处老老实实求和，再让 $x\to1^-$：

(1) $1 - 1 + 1 - 1 + \cdots$：加因子得 $\tfrac1{1+x}$（几何级数，第 22 章），令 $x\to1^-$ 得 $\tfrac12$。

(2) $1 - 2 + 3 - 4 + \cdots$：先算 $\sum_{n\ge1} n(-1)^{n-1}x^n$——对 $\tfrac{1}{1+x} = \sum(-1)^n x^n$ 两边求导再乘 $-x$，得 $\tfrac{x}{(1+x)^2}$；令 $x\to1^-$ 得 $\tfrac14$。

(3) 按黎曼的规则，$1+2+3+\cdots$（不带正负号）的延拓值是 $-\tfrac1{12}$——它需要 $\zeta$ 的完整延拓 machinery，本书记账不展开。

物理的实锤：\textbf{卡西米尔效应}（两块靠得极近的金属板之间的真空吸力）的计算里，真空能求和正是 $1+2+3+\cdots$ 的正则化，$-\tfrac1{12}$ 真实登场，且实验测得的吸力与理论吻合。\textbf{合法外推能把直觉带到多远的地方}——这道题没有答案，只有风景。
\end{shaonao}

\begin{chengshi}
本章的每个名词——解析延拓、函数方程、临界带、显式公式、随机矩阵——都只给了"观光级"解释。深挖路线图：Derbyshire《素数之恋》（最佳观光导游，零门槛）；需要公式时 Apostol《解析数论导引》；零点与随机矩阵的物理线索见《The Music of the Primes》（Sabbagh）。本书对黎曼猜想的态度与数学界一致：\textbf{参观，致敬，不装懂}。
\end{chengshi}

篇四完。数论与分析的合流，终点站是全人类的空白地图。下面换一条线路：篇五，把"弯曲"本身变成可以计算的数——先从一条高速公路的急转弯开始。
